\Needspace{4\baselineskip}
\section{\bt{Implementation}{系统实现}}
\label{sec:implementation}

\bi{\emph{Middleware and agent interface.} We implement \system{} in Rust and Python on PostgreSQL (\url{https://github.com/mmmttt000045/MAVRA}): an asynchronous HTTP server (\code{axum}, \code{tokio-postgres}, \code{deadpool-postgres}) uses \code{sqlparser} to resolve table aliases, check parameter separation, and compile SQL queries from stored definitions, exposing MCP-style JSON tools (\code{find\_metric}, snapshot-bound \code{run\_sql}, schema/column probes, \code{find\_program}, \code{find\_features}, and \code{compose\_features}).}
{\emph{中间层与智能体接口。}我们在 PostgreSQL 上用 Rust 和 Python 实现了 \system{}（\url{https://github.com/mmmttt000045/MAVRA}）：异步 HTTP 服务（\code{axum}、\code{tokio-postgres}、\code{deadpool-postgres}）借助 \code{sqlparser} 消解表别名、检查参数分离并从已存定义编译生成 SQL，向外提供 MCP 风格的 JSON 工具接口（\code{find\_metric}、绑定快照的 \code{run\_sql}、表与列探查、\code{find\_program}、\code{find\_features} 与 \code{compose\_features}）。}

\bi{\emph{Transactional version tracking and single-flight checks.} Because PostgreSQL's \code{pg\_stat\_user\_tables} counters update asynchronously outside transactions, \system{} installs statement-level \code{AFTER INSERT OR UPDATE OR DELETE OR TRUNCATE} triggers that increment a 16-shard table \code{mavra\_versions(table\_name, shard, v)} hashed by \code{pg\_backend\_pid() \% 16}, paired with \code{etl\_batch\_log} for insert-only row ranges $\Delta T$ (both hidden from agent views). Trigger updates commit with the write transaction without row-lock contention, so reading \code{SUM(v)} inside a read-only \code{REPEATABLE READ} transaction gives stamps $v_s$ matching snapshot $D_s$, while an in-memory single-flight map (\code{OnceCell}) merges concurrent checks on the same key $k(c,V)$.}
{\emph{事务级版本跟踪与并发检查合并。}由于 PostgreSQL 自带的 \code{pg\_stat\_user\_tables} 在事务外异步更新，\system{} 在表上安装语句级触发器，按 \code{pg\_backend\_pid() \% 16} 递增 16 分片版本表 \code{mavra\_versions(table\_name, shard, v)}，并配合 \code{etl\_batch\_log} 记录纯追加批次范围 $\Delta T$（两张表均对智能体隐藏）。触发器随写事务原子提交且避免行锁争用，因此只读 \code{REPEATABLE READ} 事务内读取 \code{SUM(v)} 就能拿到与快照 $D_s$ 完全对齐的版本戳 $v_s$；内存层则通过 single-flight（\code{OnceCell}）合并针对同一键 $k(c,V)$ 的并发检查。}

\bi{\emph{Lookup models.} Dense recall uses \code{BAAI/bge-m3} (top-$K=8$), and consistency scoring averages three 12-layer \code{XLM-RoBERTa} cross-encoders fine-tuned with different random seeds under fixed thresholds $(\tau_{\mathrm{low}}, \tau_{\mathrm{high}})=(0.1, 0.9)$ and cosine fallback $0.6$. On our 26-definition TPC-DS catalog, offline intent enumeration produces 98 intents expanded into 8,053 phrasings in ${\sim}500$ LLM calls; confirmed aliases take effect immediately and are included in periodic background retraining.}
{\emph{检索模型配置。}向量召回使用 \code{BAAI/bge-m3}（top-$K=8$），一致性打分取三个不同随机种子微调的 12 层 \code{XLM-RoBERTa} 交叉编码器的均值，固定阈值 $(\tau_{\mathrm{low}}, \tau_{\mathrm{high}})=(0.1, 0.9)$，余弦回退阈值 $0.6$。在包含 26 条 TPC-DS 定义的目录上，离线枚举出 98 类意图，经约 500 次大模型调用扩写为 8,053 条问法；在线确认过的别名立即生效，并在后台定期并入重训。}
